\documentclass[aspectratio=169,10pt]{beamer}
\usetheme{default}

\usepackage[T1]{fontenc}
\usepackage[utf8]{inputenc}
\usepackage{amsmath,amssymb}
\usepackage{graphicx}
\usepackage{grffile}
\usepackage{booktabs}
\usepackage{xcolor}
\usepackage{array}
\usepackage{colortbl}
\usepackage{tikz}

\graphicspath{{results/figures/}{presentation/logo/}}

\definecolor{upcblue}{HTML}{0077C8}
\definecolor{gdblue}{HTML}{2196F3}
\definecolor{gaorange}{HTML}{FF5722}
\definecolor{lightgrey}{HTML}{F5F5F5}

% ── Colors ───────────────────────────────────────────────────────────────────
\setbeamercolor{structure}{fg=upcblue}
\setbeamercolor{normal text}{fg=black,bg=white}
\setbeamercolor{frametitle}{fg=black,bg=white}
\setbeamercolor{block title}{bg=upcblue!15,fg=upcblue!80!black}
\setbeamercolor{block body}{bg=upcblue!5,fg=black}

% ── Layout ───────────────────────────────────────────────────────────────────
\setbeamersize{text margin left=0.8cm,text margin right=0.4cm}
\setbeamerfont{frametitle}{size=\normalsize,series=\bfseries}
\setbeamertemplate{navigation symbols}{}

% ── Blue left stripe on every slide ──────────────────────────────────────────
\setbeamertemplate{background}{%
  \begin{tikzpicture}
    \useasboundingbox (0,0) rectangle (\paperwidth,\paperheight);
    \fill[upcblue] (0,0) rectangle (0.45cm,\paperheight);
  \end{tikzpicture}%
}

% ── Headline: course name left + UPC logo right ───────────────────────────────
\setbeamertemplate{headline}{%
  \leavevmode
  \hbox{%
    \begin{beamercolorbox}[wd=\paperwidth,ht=1.1cm,dp=0.2cm]{headline}%
      \hspace{0.7cm}%
      \raisebox{0.3cm}{\color{gray}\footnotesize Custom Smart Adaptive Systems}%
      \hfill
      \raisebox{0.15cm}{%
        
\includegraphics[height=0.75cm]{{presentation/logo/upc-positiu-p3005 (1)}.png}%
      }%
      \hspace{0.3cm}%
    \end{beamercolorbox}%
  }%
}

% ── No footline ──────────────────────────────────────────────────────────────
\setbeamertemplate{footline}{}

% ── Frametitle: indent past the stripe ───────────────────────────────────────
\setbeamertemplate{frametitle}{%
  \vskip0.15cm
  \hspace{0.1cm}\insertframetitle
}

\title{Gradient Descent vs.\ Genetic Algorithm\\for Urban Sound Classification}
\subtitle{Custom Smart Adaptive Systems · UPC 2025/26}
\author{Leonie Greber}
\date{\today}

% ─────────────────────────────────────────────────────────────────────────────
\begin{document}

\begin{frame}
  \titlepage
\end{frame}

\begin{frame}{Outline}
  \tableofcontents
\end{frame}

% ══════════════════════════════════════════════════════════════════════════════
\section{Dataset}

\begin{frame}{Dataset: UrbanSound8K}
  \begin{columns}[T]
    \begin{column}{0.55\textwidth}
      \begin{itemize}
        \item \textbf{8\,674 audio clips}, 10 urban-sound classes
        \item Raw audio reduced to \textbf{34 numerical features} per clip
        \item Features extracted once; training operates on tabular data
      \end{itemize}
      \vspace{0.5em}
      \begin{tabular}{ll}
        \toprule
        Feature group & Dimensions \\
        \midrule
        MFCC (mel-frequency cepstral coefficients) & 13 \\
        Chroma & 12 \\
        Spectral contrast & 7 \\
        Zero-crossing rate & 1 \\
        Spectral centroid & 1 \\
        \midrule
        \textbf{Total} & \textbf{34} \\
        \bottomrule
      \end{tabular}
    \end{column}
    \begin{column}{0.42\textwidth}
      \textbf{10 classes:}\\[0.3em]
      \small
      air conditioner, car horn, children playing,\\
      dog bark, drilling, engine idling,\\
      gun shot, jackhammer, siren, street music
      \vspace{0.5em}

      \textbf{Split (stratified):}\\
      70\% train · 10\% val · 20\% test
    \end{column}
  \end{columns}
\end{frame}

\begin{frame}{Dataset: Class Distribution}
  \begin{columns}[T]
    \begin{column}{0.58\textwidth}
      \includegraphics[width=\textwidth]{class_distribution.png}
    \end{column}
    \begin{column}{0.38\textwidth}
      \begin{itemize}
        \item Most classes: $\approx$1\,000 samples
        \item \textbf{car\_horn} ($\approx$429) and \textbf{gun\_shot} ($\approx$374) are underrepresented
        \item Stratified split ensures proportional representation in every fold
        \item These two classes will show lower F1 (by construction, not model failure)
      \end{itemize}
    \end{column}
  \end{columns}
\end{frame}

\begin{frame}{Dataset: Feature Importance and Class Separability}
  \begin{columns}[T]
    \begin{column}{0.48\textwidth}
      \includegraphics[width=\textwidth]{feature_importance.png}\\
      \centering\small MFCC features dominate mutual information
    \end{column}
    \begin{column}{0.48\textwidth}
      \includegraphics[width=\textwidth]{pca_2d.png}\\
      \centering\small Classes overlap in 2D; non-linear boundaries needed
    \end{column}
  \end{columns}
  \vspace{0.2em}
  \begin{itemize}
    \item PCA retains $\approx$42\% of variance in 2 components; substantial structure remains in higher dims
    \item A non-linear model (MLP with ReLU) is therefore justified
  \end{itemize}
\end{frame}

% ══════════════════════════════════════════════════════════════════════════════
\section{Architecture}

\begin{frame}{MLP Architecture}
  \begin{columns}[T]
    \begin{column}{0.52\textwidth}
      \textbf{Forward pass:}
      \[
        \mathbf{z}^{(1)} = W_1\mathbf{x} + \mathbf{b}_1, \quad
        \mathbf{h} = \text{ReLU}(\mathbf{z}^{(1)})
      \]
      \[
        \mathbf{z}^{(2)} = W_2\mathbf{h} + \mathbf{b}_2 \quad\text{(logits)}
      \]
      \vspace{0.3em}
      {\scriptsize\texttt{Input(34) $\to$ Linear(34$\to$20) $\to$ ReLU $\to$ Linear(20$\to$10)}}

      \vspace{0.7em}
      \textbf{Parameter count:}\\[0.3em]
      \begin{tabular}{lcc}
        \toprule
        Layer & Formula & Count \\
        \midrule
        Linear(34$\to$20) & $34\times20 + 20$ & 700 \\
        Linear(20$\to$10) & $20\times10 + 10$ & 210 \\
        \midrule
        \textbf{Total} & & \textbf{910} \\
        \bottomrule
      \end{tabular}
    \end{column}
    \begin{column}{0.44\textwidth}
      \begin{itemize}
        \item \textbf{Loss:} CrossEntropyLoss (log-softmax + NLL)
        \item Forward returns \emph{raw logits}; no explicit softmax layer
        \item \textbf{Same architecture} for both GD and GA; only the weight-update mechanism differs
        \item GA treats the MLP as a black box: no gradients required
      \end{itemize}
    \end{column}
  \end{columns}
\end{frame}

\begin{frame}{Why Hidden Size = 20? (Bias-Variance Constraint)}
  \begin{columns}[T]
    \begin{column}{0.54\textwidth}
      \begin{block}{Rule of thumb: 5--10 training samples per parameter}
        $\dfrac{6{,}072\text{ train samples}}{910\text{ parameters}} \approx 6.7
        \;\Rightarrow$ within the safe zone
      \end{block}
      \vspace{0.4em}
      \begin{tabular}{cccc}
        \toprule
        Hidden & Params & Samples/param & Risk \\
        \midrule
        10 & 460 & 13.2 & underfitting \\
        \rowcolor{gdblue!20}
        \textbf{20} & \textbf{910} & \textbf{6.7} & \textbf{good} \\
        30 & 1{,}360 & 4.5 & overfitting \\
        40 & 1{,}810 & 3.4 & overfitting \\
        50 & 2{,}260 & 2.7 & overfitting \\
        \bottomrule
      \end{tabular}
    \end{column}
    \begin{column}{0.42\textwidth}
      \begin{itemize}
        \item Larger hidden layer can fit train data better but \emph{generalises worse} on a dataset of this size
        \item This choice is \emph{confirmed empirically} by the parameter sweep in Section~4 (train$-$val gap widens for $h > 20$)
        \item Weight decay ($\lambda=10^{-4}$) adds L2 regularisation as a second line of defence
      \end{itemize}
    \end{column}
  \end{columns}
\end{frame}

% ══════════════════════════════════════════════════════════════════════════════
\section{Mathematical Background}

\begin{frame}{GD: Loss Function and Gradient}
  \textbf{Cross-entropy loss} (multi-class, $N$ samples, $C=10$ classes):
  \begin{align}
    \mathcal{L}(\theta) &= -\frac{1}{N}\sum_{i=1}^{N} \log \hat{p}_{y_i}(\mathbf{x}_i;\,\theta)
    \tag{1}\\[4pt]
    \hat{p}_k(\mathbf{x}) &= \frac{e^{z_k(\mathbf{x})}}{\displaystyle\sum_{j=1}^{C} e^{z_j(\mathbf{x})}}
    \quad\text{(softmax)}
    \tag{2}
  \end{align}

  \textbf{Vanilla gradient descent update:}
  \begin{equation}
    \theta_{t+1} = \theta_t - \eta\,\nabla_\theta \mathcal{L}(\theta_t)
    \tag{3}
  \end{equation}

  \begin{itemize}
    \item $\nabla_\theta \mathcal{L}$ is computed via \textbf{backpropagation} (chain rule through both linear layers)
    \item Mini-batch size 64 — trade-off between gradient noise (regularisation) and convergence speed
    \item With weight decay $\lambda$: effective loss is $\mathcal{L}(\theta) + \tfrac{\lambda}{2}\|\theta\|^2$
  \end{itemize}
\end{frame}

\begin{frame}{GD: Adam Optimizer}
  Adam maintains \emph{per-parameter} running estimates of gradient mean and variance:
  \begin{align}
    m_t &= \beta_1 m_{t-1} + (1-\beta_1)\,g_t \tag{4}\\
    v_t &= \beta_2 v_{t-1} + (1-\beta_2)\,g_t^2 \tag{5}\\
    \hat{m}_t &= \frac{m_t}{1-\beta_1^t}, \quad
    \hat{v}_t = \frac{v_t}{1-\beta_2^t}
    \quad\text{(bias correction)} \tag{6}\\
    \theta_{t+1} &= \theta_t - \frac{\eta}{\sqrt{\hat{v}_t}+\epsilon}\,\hat{m}_t
    \tag{7}
  \end{align}

  \begin{columns}[T]
    \begin{column}{0.45\textwidth}
      \textbf{Parameters used:}\\[0.2em]
      $\eta = 10^{-3},\quad \beta_1 = 0.9,\quad \beta_2 = 0.999$\\
      $\epsilon = 10^{-8},\quad \lambda_{\text{wd}} = 10^{-4}$
    \end{column}
    \begin{column}{0.50\textwidth}
      \begin{itemize}
        \item $\sqrt{\hat{v}_t}$ normalises the step: large-gradient params take smaller steps
        \item Bias correction avoids small initial estimates biasing early updates
      \end{itemize}
    \end{column}
  \end{columns}
\end{frame}

\begin{frame}{GD: Training Dynamics}
  \begin{center}
    \includegraphics[width=0.85\textwidth]{gd_curves.png}
  \end{center}
  \begin{itemize}
    \item Loss (log scale) falls near-linearly: confirms exponential decay, plateauing after epoch~80
    \item Train and val accuracy track closely, confirming hidden$=20$ does not overfit
    \item Log scale makes the convergence rate visible; the slope flattens after epoch~80
  \end{itemize}
\end{frame}

\begin{frame}{GA: Operators}
  The GA searches the 910-dimensional weight space \emph{without gradients}.

  \begin{columns}[T]
    \begin{column}{0.48\textwidth}
      \textbf{Fitness (improved):}
      \begin{equation}
        f(\theta) = -\mathcal{L}_{\text{val}}(\theta)
        \tag{8}
      \end{equation}
      Higher fitness = lower validation loss; continuous signal vs.\ coarse accuracy.

      \vspace{0.5em}
      \textbf{BLX-$\alpha$ crossover} ($\alpha=0.5$):
      \begin{equation}
        o_i \sim \mathcal{U}[l_i,\,u_i]
        \tag{9}
      \end{equation}
      \small$l_i = \min(p_i,q_i) - \alpha|p_i-q_i|$,\;
      $u_i = \max(p_i,q_i) + \alpha|p_i-q_i|$\\[0.2em]
      \normalsize Allows offspring outside parent range.
    \end{column}
    \begin{column}{0.48\textwidth}
      \textbf{Gaussian mutation} ($p_m = 0.05$ per gene):
      \begin{equation}
        \theta_i \leftarrow \theta_i + \mathcal{N}(0,\,\sigma^2),\;\sigma=0.05
        \tag{10}
      \end{equation}
      Per-gene probability: $\approx$46 of 910 weights perturbed per individual.

      \vspace{0.5em}
      \textbf{Selection:} Tournament of size 5: sample 5 individuals, keep the best.
      Repeat $N-k$ times to fill the next generation.
    \end{column}
  \end{columns}
\end{frame}

\begin{frame}{GA: Algorithm per Generation}
  \begin{block}{One GA generation (pseudocode)}
    \begin{enumerate}
      \item Evaluate fitness for all individuals with invalid fitness
      \item Select $N - k$ offspring via tournament selection
      \item For each consecutive pair: apply BLX-$\alpha$ crossover with probability $c = 0.7$
      \item For each offspring: apply Gaussian mutation with probability $p_m = 0.05$
      \item Preserve top-$k$ \textbf{elite} individuals unchanged (elitism, $k=5$)
      \item Replace population: offspring $\cup$ elites ($N$ individuals total)
      \item Update Hall of Fame (best individual ever seen across all generations)
    \end{enumerate}
  \end{block}

  \begin{itemize}
    \item \textbf{Baseline GA} used raw accuracy fitness and \texttt{eaSimple} (no elitism); test acc 0.377
    \item \textbf{Key improvements:} fitness $= -\mathcal{L}_{\text{val}}$ + elitism $\Rightarrow$ test acc 0.621
    \item \textbf{Early stopping:} patience of 50 generations without improvement; avoids wasted compute
  \end{itemize}
\end{frame}

\begin{frame}{GA: Convergence (Improved)}
  \begin{center}
    \includegraphics[width=0.85\textwidth]{ga_convergence.png}
  \end{center}
  \begin{itemize}
    \item Left: validation loss of best individual (HoF) vs population average; gap shows elites pull ahead
    \item Right: accuracy of improved GA far exceeds baseline across all generations
    \item Population std (shaded) gradually narrows; diversity decreases as elites dominate
  \end{itemize}
\end{frame}

% ══════════════════════════════════════════════════════════════════════════════
\section{Parameter Tuning}

\begin{frame}{GD: Hyperparameter Sweeps (Learning Rate and Weight Decay)}
  \begin{columns}[T]
    \begin{column}{0.49\textwidth}
      \centering
      \includegraphics[width=\textwidth]{tuning_gd_lr.png}\\[0.2em]
      \small Chosen: $\eta = 10^{-3}$
    \end{column}
    \begin{column}{0.49\textwidth}
      \centering
      \includegraphics[width=\textwidth]{tuning_gd_wd.png}\\[0.2em]
      \small Chosen: $\lambda = 10^{-4}$
    \end{column}
  \end{columns}
  \vspace{0.4em}
  \small All sweeps fix other parameters at chosen values; 150 epochs; seed fixed.
  Hidden-size sweep: see next slide.
\end{frame}

\begin{frame}{Bias-Variance Tradeoff: Hidden Size in Detail}
  \begin{center}
    \includegraphics[width=0.82\textwidth]{tuning_gd_hidden.png}
  \end{center}
  \begin{columns}[T]
    \begin{column}{0.48\textwidth}
      \begin{itemize}
        \item $h=10$: both train and val acc are modest; \textbf{underfitting} (insufficient capacity)
        \item $h=20$: small train$-$val gap, good val acc; \textbf{well-calibrated}
      \end{itemize}
    \end{column}
    \begin{column}{0.48\textwidth}
      \begin{itemize}
        \item $h=30$--$50$: train acc rises, val acc plateaus/drops; \textbf{overfitting}, gap widens
        \item Weight decay compresses the gap but also limits peak accuracy at high $h$
      \end{itemize}
    \end{column}
  \end{columns}
\end{frame}

\begin{frame}{GA: Hyperparameter Sweeps}
  \begin{columns}[T]
    \begin{column}{0.33\textwidth}
      \centering
      \includegraphics[width=\textwidth]{tuning_ga_pop.png}\\
      \small Chosen: $N = 100$
    \end{column}
    \begin{column}{0.33\textwidth}
      \centering
      \includegraphics[width=\textwidth]{tuning_ga_mutpb.png}\\
      \small Chosen: $p_m = 0.05$
    \end{column}
    \begin{column}{0.33\textwidth}
      \centering
      \includegraphics[width=\textwidth]{tuning_ga_elite.png}\\
      \small Chosen: $k = 5$
    \end{column}
  \end{columns}
  \vspace{0.3em}
  \small Sweeps use 100 generations each; $c = 0.7$ (crossover prob) fixed. Hatched bar = chosen value.

  \vspace{0.3em}
  \begin{itemize}
    \item Larger populations explore more but with diminishing returns and higher per-generation cost
    \item Low mutation ($p_m = 0.01$) causes premature convergence; high ($p_m = 0.2$) destroys good solutions
    \item Elitism $k=5$ reliably outperforms $k=1$: a single elite can be lost to crossover; five provide a stable backbone
  \end{itemize}
\end{frame}

% ══════════════════════════════════════════════════════════════════════════════
\section{Results}

\begin{frame}{Results: Accuracy and Training Time}
  \begin{columns}[T]
    \begin{column}{0.52\textwidth}
      \includegraphics[width=\textwidth]{comparison_summary.png}
    \end{column}
    \begin{column}{0.44\textwidth}
      \begin{tabular}{lcc}
        \toprule
        & \textcolor{gdblue}{\textbf{GD}} & \textcolor{gaorange}{\textbf{GA}} \\
        \midrule
        Test accuracy & \textbf{0.811} & 0.621 \\
        Training time & $\approx$37\,s & $\approx$124\,s \\
        Macro F1 & \textbf{$\approx$0.81} & $\approx$0.60 \\
        Stopped at & 150 epochs & gen $\approx$799 \\
        \bottomrule
      \end{tabular}
      \vspace{0.5em}
      \begin{itemize}
        \item GD is $\approx$30\% more accurate and $3.4\times$ faster
        \item GA's gradient-free search struggles in 910 dimensions
        \item Elitism + $-\mathcal{L}$ fitness roughly doubled GA accuracy over baseline (0.377 $\to$ 0.621)
      \end{itemize}
    \end{column}
  \end{columns}
\end{frame}

\begin{frame}{Results: Convergence}
  \begin{center}
    \includegraphics[width=\textwidth]{comparison_convergence.png}
  \end{center}
  \begin{columns}[T]
    \begin{column}{0.50\textwidth}
      \begin{itemize}
        \item GD converges in $<50$ epochs; val accuracy plateaus near $81\%$
        \item Rapid early gains reflect Adam's adaptive step-size
      \end{itemize}
    \end{column}
    \begin{column}{0.48\textwidth}
      \begin{itemize}
        \item GA improves slowly across ${\approx}800$ generations; best val accuracy ${\approx}62\%$
        \item 1 GA generation $= 100$ individuals $\times$ val passes; $\gg10\times$ more compute per point
      \end{itemize}
    \end{column}
  \end{columns}
\end{frame}

\begin{frame}{Results: Per-class F1 Score}
  \begin{center}
    \includegraphics[width=0.90\textwidth]{comparison_f1.png}
  \end{center}
  \begin{itemize}
    \item GD dominates on all 10 classes; both methods struggle most with \textit{gun\_shot} and \textit{car\_horn}
    \item Low F1 on those two classes reflects the class imbalance, not model failure
  \end{itemize}
\end{frame}

\begin{frame}{Results: Confusion Matrices}
  \begin{center}
    \includegraphics[width=\textwidth]{comparison_confusion.png}
  \end{center}
  \begin{itemize}
    \item GA confusion matrix shows more off-diagonal mass; less discriminative across all classes
    \item GD confuses acoustically similar pairs (e.g.\ jackhammer/drilling); GA confuses almost everything
  \end{itemize}
\end{frame}

\begin{frame}{Conclusion}
  \begin{columns}[T]
    \begin{column}{0.48\textwidth}
      \textbf{Key findings:}
      \begin{itemize}
        \item GD (Adam) clearly dominates: 81.1\% vs 62.1\% accuracy, $3\times$ faster
        \item Continuous $-\mathcal{L}$ fitness + elitism are \emph{essential} for GA (baseline: 37.7\%)
        \item Hidden $= 20$ is empirically confirmed as the optimal capacity for this dataset
        \item Both methods agree: \textit{gun\_shot} and \textit{car\_horn} are hardest (class imbalance)
      \end{itemize}
    \end{column}
    \begin{column}{0.48\textwidth}
      \textbf{When does GA make sense?}
      \begin{itemize}
        \item Non-differentiable objectives (hardware-in-the-loop, discrete search spaces)
        \item Black-box optimisation where gradient computation is impossible
      \end{itemize}
      \vspace{0.5em}
      \textbf{Future work:}
      \begin{itemize}
        \item CMA-ES: adapts the mutation covariance matrix, far more efficient in high dimensions
        \item Hybrid: GD warm-start $\to$ GA fine-tune for non-smooth loss surfaces
        \item Neuroevolution of architectures (NEAT): evolve topology, not just weights
      \end{itemize}
    \end{column}
  \end{columns}
\end{frame}

\end{document}
